\begin{figure}[h]
	\centering
	\begin{tikzpicture}
		\coordinate (O) at (0,0);
		\coordinate (A) at ({1.25*cos(90+60)}, {1.25*sin(90+60)});
		
		\draw (0,-1.75) -- (0,2) 
			node[right, yshift=1.5]{$y$};
		
		\draw[blue!80] (O) circle [x radius=1.25cm, y radius=1.25cm];
		\draw[rotate around={-50:(O)}, red!80] (O) circle [x radius=1.25cm, y radius=2.5cm];
		
		\draw (-2.25,0) -- (2.25,0) 
			node[below right]{$x$};
		\draw[dashed, >=latex, <->, rotate around={-50:(O)}] (O) --+ (0,2.5) 
			node[midway, fill=white, inner sep=0.1pt] {\footnotesize$\lambda_{max}$};
		\draw[dashed, >=latex, <->, rotate around={-130:(O)}] (O) --+ (0,1.25) 
			node[midway, fill=white, inner sep=0.1pt] {\footnotesize$\lambda_{min}$};
		
		\draw[-{Latex[length=2mm]}] (0,0) + (0:2cm) arc (0:40:2cm) 
			node[midway, xshift=-2mm]{$\phi$};
			
		%\draw[>=latex, ->, rotate around={60:(O)}] (O) --+ (0,1.25);
		\draw[>=latex, ->, blue] (O) -- (A)
			node[midway, above]{$\delta_x$};
		\draw[dashed, gray, rotate around={-50:(A)}] (A) --+ (0,2);
		\draw[>=latex, ->, red] ($(A) + ({1.25*cos(90-50)},{1.25*sin(90-50)})$) -- (O)
			node[midway, right]{$\delta_f$};
	\end{tikzpicture}
	\caption{Exemple d'ellipse avec un déplacement $\delta_x$ (bleu) sur le cercle unitaire, et la force de restitution (rouge) non-colinéaire induite, que l'on obtient grâce à l'intersection de la droite parallèle à l'axe principal passant par $\delta_x$ et l'ellipse.}
	\label{fig_ellipse_dolan}
\end{figure}